\documentclass[12pt,a4paper]{article}
\usepackage[utf8]{inputenc}
\usepackage{amsmath,amssymb,amsfonts}
\usepackage{graphicx}
\usepackage{hyperref}
\usepackage{tikz}
\usetikzlibrary{shapes.geometric, arrows.meta, positioning, fit, backgrounds}
\usepackage{geometry}
\usepackage{tcolorbox}
\usepackage{listings}
\usepackage{xcolor}
\usepackage{titlesec}

\geometry{a4paper, margin=1in}

% --- Code Block Styling ---
\definecolor{codegreen}{rgb}{0,0.6,0}
\definecolor{codegray}{rgb}{0.5,0.5,0.5}
\definecolor{codepurple}{rgb}{0.58,0,0.82}
\definecolor{backcolour}{rgb}{0.95,0.95,0.92}

\lstdefinestyle{mystyle}{
    backgroundcolor=\color{backcolour},   
    commentstyle=\color{codegreen},
    keywordstyle=\color{magenta},
    numberstyle=\tiny\color{codegray},
    stringstyle=\color{codepurple},
    basicstyle=\ttfamily\footnotesize,
    breakatwhitespace=false,         
    breaklines=true,                 
    captionpos=b,                    
    keepspaces=true,                 
    numbers=left,                    
    numbersep=5pt,                  
    showspaces=false,                
    showstringspaces=false,
    showtabs=false,                  
    tabsize=2
}
\lstset{style=mystyle}

% --- Headers Styling ---
\titleformat{\section}{\Large\bfseries\color{blue!60!black}}{\thesection}{1em}{}
\titleformat{\subsection}{\large\bfseries\color{blue!40!black}}{\thesubsection}{1em}{}

\title{\Huge \textbf{The Complete Guide to CodeAudit AI} \\ 
\Large From Foundational Deep Learning to Multi-Agent Vector Retrieval}
\author{Comprehensive Technical Reference}
\date{\today}

\begin{document}

\maketitle
\tableofcontents
\newpage

\section{Introduction}
This document serves as the definitive guide to understanding and replicating the \textbf{CodeAudit AI} system. It is designed to take a reader with zero prior knowledge of Deep Learning and guide them all the way to building a complex Multi-Agent Retrieval-Augmented Generation (RAG) architecture. Whether you are defending the project to a review panel or teaching a peer how to build their own system, this reference contains the complete theoretical foundation and technical implementation details.

% =========================================================
\section{Part 1: The Foundations of Modern AI}
% =========================================================

\subsection{Machine Learning vs. Deep Learning}
\begin{itemize}
    \item \textbf{Machine Learning (ML):} Algorithms that learn patterns from data instead of being explicitly programmed. Example: A spam filter learning that emails with the word "Lottery" are usually spam.
    \item \textbf{Deep Learning (DL):} A subset of ML that uses \textit{Artificial Neural Networks} with many layers ("deep" layers) to learn highly complex patterns.
\end{itemize}

\subsection{Large Language Models (LLMs) and Transformers}
\begin{itemize}
    \item \textbf{Transformers:} The neural network architecture (invented in 2017) that makes LLMs possible. Transformers process an entire sequence of text simultaneously and learn the "attention" or relationship between words.
    \item \textbf{How LLMs work:} Fundamentally, an LLM takes a sequence of words (the prompt) and mathematically predicts the most likely \textit{next word}.
\end{itemize}

\begin{tcolorbox}[colback=red!5!white,colframe=red!75!black,title=The Context Window Problem]
\textbf{Why don't we just ask an LLM to read the entire GitHub repository directly?} \\
An LLM has a strict memory limit called a \textbf{Context Window} (e.g., 8,000 to 128,000 tokens). A candidate's GitHub might contain millions of tokens of code. If you feed it all to the LLM, it crashes. If you ask it to judge a skill without feeding it the code, it will \textbf{hallucinate} (guess facts). We must search and filter the code \textit{first}.
\end{tcolorbox}

% =========================================================
\section{Part 2: Vector Embeddings and Semantic Search}
% =========================================================

\subsection{The Math Behind Meaning}
You cannot use standard text searching ("Ctrl+F") to verify skills in code. Searching the text "Machine Learning" will fail if the developer wrote \texttt{import tensorflow as tf}. We need \textbf{Semantic Search}.

We solve this using \textbf{Vector Embeddings}. An embedding model converts text or code into a coordinate in a high-dimensional mathematical space.
Imagine a simplified 3D coordinate system $[X, Y, Z]$ representing $[Animal, Vehicle, Code]$:
\begin{itemize}
    \item "Dog" $\rightarrow [0.9, 0.1, 0.0]$
    \item "Car" $\rightarrow [0.1, 0.9, 0.0]$
    \item \texttt{print("Hello")} $\rightarrow [0.0, 0.1, 0.9]$
\end{itemize}

In our system, we use the \texttt{all-MiniLM-L6-v2} model. It converts any snippet of code into an array of \textbf{384 numbers} (384 dimensions).

\subsection{Cosine Similarity}
To find the closest match between a resume claim ("Deep Learning") and thousands of files of code, we calculate the \textbf{Cosine Similarity}. This measures the \textit{angle} between two 384-dimensional vectors. An angle of 0 degrees means a perfect semantic match (Cosine = 1.0).

\begin{lstlisting}[language=Python, caption=Calculating Semantic Similarity]
from sentence_transformers import SentenceTransformer
from sklearn.metrics.pairwise import cosine_similarity

model = SentenceTransformer('all-MiniLM-L6-v2')

claim_vector = model.encode(["Machine Learning"])
code_vector = model.encode(["import torch.nn as nn"])

# The text differs completely, but the math proves they are semantically similar!
similarity = cosine_similarity(claim_vector, code_vector)
print(similarity) # Output: > 0.65
\end{lstlisting}

% =========================================================
\section{Part 3: Retrieval-Augmented Generation (RAG) \& ChromaDB}
% =========================================================

\subsection{The RAG Pipeline}
RAG prevents LLM hallucinations through a two-step process:
\begin{enumerate}
    \item \textbf{Retrieval:} Search our database for the exact code snippets relevant to the candidate's skill claim.
    \item \textbf{Generation:} Hand those specific code snippets to the LLM (like giving an open book to a student) and ask it to verify the claim based \textit{only} on the provided code.
\end{enumerate}

\subsection{Vector Databases (ChromaDB)}
You cannot calculate Cosine Similarity manually against 10,000 code files in a standard SQL database; it is computationally impossible to do efficiently. We use a \textbf{Vector Database} (ChromaDB) which organizes vectors using advanced HNSW (Hierarchical Navigable Small World) graph algorithms for millisecond retrieval.

\begin{lstlisting}[language=Python, caption=Storing and Querying vectors in ChromaDB]
import chromadb
client = chromadb.Client()
collection = client.create_collection("github_code")

# 1. Store the embedded code chunks
collection.add(
    documents=["import pandas as pd", "def build_web_app(): pass"],
    metadatas=[{"user": "candidate123"}, {"user": "candidate123"}],
    ids=["chunk1", "chunk2"]
)

# 2. Retrieve the top match
results = collection.query(
    query_texts=["Data Analysis"],
    n_results=1
)
# Instantly returns: "import pandas as pd"
\end{lstlisting}

% =========================================================
\section{Part 4: The Multi-Agent Architecture}
% =========================================================

An "Agent" is an LLM wrapped in specific programming logic that gives it a defined persona, specific tools, and a distinct task. Breaking a monolithic task down into specialized agents increases accuracy and speed.

\subsection{Phase 1: Extraction Agent}
\textbf{Goal:} Convert unstructured PDF data into structured JSON. \\
\textbf{Logic:} We utilize \texttt{PyMuPDF} to extract text, then prompt an LLM to parse out the GitHub URL and a normalized list of technical skills.

\subsection{Phase 2: Ingestion Agent (Crucial Step)}
\textbf{Goal:} Download candidate code via the GitHub REST API and prepare it for embedding. \\
\textbf{Chunking Logic:} If you embed an entire 10,000-line file, the resulting vector becomes "diluted" and loses specific meaning. We must slice the code into \textbf{800-character chunks}. \\
\textbf{Sliding Window Overlap:} If we hard-cut every 800 characters, we might cut a critical function in half. We use a \textbf{100-character sliding overlap}, ensuring the end of Chunk 1 is duplicated at the beginning of Chunk 2.

\subsection{Phase 3: Indexing Agent}
\textbf{Goal:} Convert chunks to vectors and store them in ChromaDB. \\
\textbf{Multi-Tenancy Rule:} Every chunk stored in ChromaDB must contain metadata tagging the candidate's username (e.g., \texttt{\{"user": "johndoe"\}}). Without this, the query engine might accidentally verify Candidate A using Candidate B's code.

\subsection{Phase 4: LLM Judge Agent (Chain-of-Thought)}
\textbf{Goal:} Analyze the retrieved code chunks and output a final decision. \\
\textbf{Chain-of-Thought (CoT):} We force the LLM to "think out loud" step-by-step before deciding. This drastically reduces false positives.

\begin{tcolorbox}[colback=blue!5!white,colframe=blue!75!black,title=The System Prompt (Chain-of-Thought)]
\textbf{Role:} You are a strict technical recruiter. \\
\textbf{Skill Claim:} \{skill\} \\
\textbf{Evidence:} \{retrieved\_chunks\} \\
\vspace{0.5em}
\textbf{Task Requirements:} \\
Step 1: Explicitly state whether the code imports libraries related to the \{skill\}. \\
Step 2: Explicitly state whether the code executes actual functions/methods related to the \{skill\}, or if it is just a text comment. \\
Step 3: Based on Steps 1 and 2, output your final decision: "Verified" or "Unsubstantiated".
\end{tcolorbox}

% =========================================================
\section{Defense Q\&A}
% =========================================================
\begin{itemize}
    \item \textbf{"Did you just wrap the OpenAI API?"} \newline \textit{No. We built a custom asynchronous pipeline handling AST-aware code chunking, dense vector indexing, and Agentic orchestration. The LLM is only utilized as the final semantic router and judge.}
    \item \textbf{"Why not fine-tune a model instead of using RAG?"} \newline \textit{Fine-tuning teaches an LLM a new format or style, but it is highly inefficient for storing factual knowledge. Developers push new code to GitHub daily. RAG allows our system to query real-time, ground-truth data without needing to retrain neural networks.}
    \item \textbf{"How do you prevent False Positives?"} \newline \textit{By utilizing a strict Chain-of-Thought rubric with a zero-temperature LLM. We force the LLM to identify explicit Abstract Syntax Tree (AST) execution patterns—if a skill is only mentioned in a string or comment, it is rejected.}
\end{itemize}

\end{document}
